\scanpage{7}
Nous allons étudier de plus près le morphisme
\[
\underline{\mathrm{Pic}}_{A/S}\longrightarrow\underline{\mathrm{Hom}}_{S\text{-gr}}(A,A')\tag*{$(*)$}
\]
\margin{hom pointé, \deleted{\unsure{ou|de|$\leftarrow$}} hom de groupes}
i.e.
\[
\underline{\mathrm{Pic}}_{A/S}\times_S A\longrightarrow A'\tag*{$(**)$}
\]
défini par
\[
(\xi\unsure{,|;}t)\rightsquigarrow t_*(\xi)-\xi
\]

\noindent\llap{\smash{\rule[-3.6ex]{0.4pt}{4ex}}\quad}\uline{Théorème} L'hom $(*)$ est un hom de $S$-Groupes \unsure{?|$\xi$|:} (i.e. $(**)$ est un accouplement), et son noyau est \uline{exactement} $A'$, \deleted{\unsure{la}} \inserted{et même :} \deleted{\unsure{dès que|pour que} \unsure{l'on a|une}} une section $\xi$ de $\underline{\mathrm{Pic}}_{A/S}$ sur $S$ est une \deleted{\unsure{est nulle|distincte} \unsure{la|de} \unsure{loi}} section de $A'$ \unsure{ssi|ses} on a
\[
\boxed{\pi^*(\xi)=\mathrm{pr}_1^*(\xi)+\mathrm{pr}_2^*(\xi)}
\]
où $\pi:A\times_{\unsure{S|s}} A\longrightarrow A$ est la \unsure{loi de} composition.

\noindent\llap{\smash{\rule[-1.8ex]{0.4pt}{2ex}}\quad}\uline{Corollaire} On déduit de $(*)$ par passage au quotient une \uline{immersion fermée}
\[
\underline{\text{\unsure{NérSév|NérSé}}}_{A/S}\longrightarrow\underline{\mathrm{Hom}}_S(A,A')\tag*{$(**)$}
\]
% slanted
\margin{$\underline{\text{NérSév}}_{A/S}=\underline{\mathrm{Pic}}_{A/S}/\underline{\mathrm{Pic}}^0_{A/S}$ par définition}
compatible avec les structures de groupe.

Démonstration du théorème. L'additivité en $t$ est immédiate grâce au fait que le morphisme $A\unsure{\text{---}|\longrightarrow}A'$ défini par $s$ transforme section unité en section unité. L'additivité en $s$ est de toute façon triviale. Le morphisme $A'\times A\longrightarrow A'$ induit par $(**)$
% archive sheet number: 6

